% P05.1 one-dimensional regularity convention repair, 25 September 2026.

Write $\mathrm{PA}_{\mathbf Z}(\mathbf X)$ for the literal class in
\cite[Definition~2.28]{BMPS26}: the linear part on every face is integral and
the constant term is rational.  Write
$\mathrm{PA}_{\mathbf R}(\mathbf X)$ for its real scalar extension, consisting
of continuous functions which are affine with arbitrary real coefficients on
the faces of a finite rational polyhedral structure.  Write
$\mathrm{PAPC}_{\mathbf Z}$ and
$\mathrm{PC}_{\mathrm{reg},\mathbf Z}$ for the source's literal classes, and
define
$\mathrm{PAPC}_{\mathbf R}(\mathbf X,\gamma)$ and
$\mathrm{PC}_{\mathrm{reg},\mathbf R}(\mathbf X,\gamma)$ from
$\mathrm{PA}_{\mathbf R}$ exactly as the source defines $\mathrm{PAPC}$ and
$\mathrm{PC}_{\mathrm{reg}}$ from $\mathrm{PA}_{\mathbf Z}$.  Thus the latter
space consists of uniform limits of decreasing sequences in
$\mathrm{PAPC}_{\mathbf R}(\mathbf X,\gamma)$.  We say that
$(\mathbf X,\gamma)$ has real-polyhedral regularity if
$P_\gamma(u)\in\mathrm{PC}_{\mathrm{reg},\mathbf R}(\mathbf X,\gamma)$ for
every bounded $u\in\mathrm{PA}_{\mathbf R}(\mathbf X)$.

\begin{theorem}
\label{thm:frontier-polyhedral-one-dimensional-real-regularity}
Let $\mathbf X$ be a one-dimensional balanced rational polyhedral space
admitting a finite structure, and let
$\gamma\in\mathrm{PA}_{\mathbf R}(\mathbf X)\cap
\mathrm{PC}(\mathbf X)$ be positive at infinity.  If
$u\in\mathrm{PA}_{\mathbf R,b}(\mathbf X)$ and $\Pi$ is a common finite
rational polyhedral structure on which $\gamma$ and $u$ are affine, then
\[
 P_\gamma(u)\in
 \mathrm{PAPC}_{\mathbf R}(\mathbf X,\gamma)
\]
and $P_\gamma(u)$ is affine on every edge of $\Pi$.  In particular,
$(\mathbf X,\gamma)$ has real-polyhedral regularity.
\end{theorem}

\begin{proposition}
\label{prop:frontier-polyhedral-integral-regularity-counterexample}
For the literal integral-slope class, regularity can fail even on the complete
affine line.  More precisely, let $\mathbf X=\mathbf R$,
$\gamma(x)=|x|$, and
\[
 u(x)=
 \begin{cases}
  0,&x\leq1,\\
  1-x,&1\leq x\leq2,\\
  -1,&x\geq2.
 \end{cases}
\]
Then $\gamma\in\mathrm{PAPC}_{\mathbf Z}(\mathbf R)$ is positive at infinity,
$u\in\mathrm{PA}_{\mathbf Z,b}(\mathbf R)$, and
\[
 P_\gamma(u)(x)=
 \begin{cases}
  -x,&x\leq0,\\
  x/2,&0\leq x\leq2,\\
  x-1,&x\geq2.
 \end{cases}
\]
Moreover every
$\psi\in\mathrm{PAPC}_{\mathbf Z}(\mathbf R,\gamma)$ satisfies
\[
 \lVert\psi-P_\gamma(u)\rVert_{L^\infty(\mathbf R)}\geq\frac14.
\]
Consequently
$P_\gamma(u)\notin\mathrm{PC}_{\mathrm{reg},\mathbf Z}
(\mathbf R,\gamma)$, and the literal integral-slope pair
$(\mathbf R,|\cdot|)$ does not have the regularity property of
\cite[Theorem~6.4 and Definition~6.5]{BMPS26}.
\end{proposition}

Thus the one-dimensional assertion in
\cite[Proposition~6.13]{BMPS26} is valid after real scalar extension but is
incompatible, as written, with the integral-slope convention in the source's
Definition~2.28.  The theorem and proposition make no higher-dimensional
regularity claim.
